\documentclass[conference]{IEEEtran}
\IEEEoverridecommandlockouts

\usepackage{cite}
\usepackage{amsmath,amssymb,amsfonts}
\usepackage{graphicx}
\usepackage{textcomp}
\usepackage{xcolor}
\usepackage{booktabs}
\usepackage{multirow}
\usepackage{listings}
\usepackage{url}

\def\BibTeX{{\rm B\kern-.05em{\sc i\kern-.025em b}\kern-.08em
    T\kern-.1667em\lower.7ex\hbox{E}\kern-.125emX}}

% compact monospace listings for command examples
\lstset{
  basicstyle=\ttfamily\footnotesize,
  breaklines=true,
  columns=fullflexible,
  keepspaces=true,
  frame=none,
  aboveskip=2pt,
  belowskip=2pt,
}

\begin{document}

\title{nosh: A Compact Fine-Tuned Language Model for\\
Natural-Language-to-Bash Command Generation\\
on Consumer Hardware}

\author{\IEEEauthorblockN{Author Name}
\IEEEauthorblockA{\textit{Independent Research} \\
\textit{}\\
confusedstudent13@gmail.com}
}

\maketitle

\begin{abstract}
Command-line interfaces remain a barrier for non-expert users, yet deploying
large language models (LLMs) to translate natural language into shell commands
is impractical on personal machines. We present \textbf{nosh}, a 1.5-billion
parameter model that converts natural-language requests into executable Bash
commands and runs comfortably on a commodity laptop with no GPU. Starting from
Qwen2.5-Coder-1.5B-Instruct, we apply parameter-efficient LoRA fine-tuning and
show that the dominant factor in command-generation quality is the
\emph{difficulty distribution} of the training data rather than the fine-tuning
hyperparameters: an initial model trained on short, tool-homogeneous corpora is
statistically indistinguishable from the base model on complex commands
($4.18$ vs. $3.50$ mean judge score is the gap we later close), whereas adding a
corpus of long, multi-tool pipelines raises the mean LLM-judge score on a
held-out complex set from $3.50$ to $4.18$ (out of $5$) and nearly doubles
exact-match accuracy from $0.26$ to $0.46$. We evaluate with a reference-aware
LLM-as-judge rubric that scores tool choice, flag coverage, argument fidelity,
constraint satisfaction, and safety, complemented by deterministic syntax and
formatting checks. Finally, we conduct a deployment study over four quantization
tiers and four few-shot prompt lengths, finding that a 4-bit ($Q4\_K\_M$,
$940$\,MB) model with a two-example prompt is the quality/footprint sweet spot,
sustaining a peak resident memory of $\approx 1.7$\,GB and a cold-start latency
of $1.4$--$1.6$\,s per query on CPU. The resulting tool integrates with the
shell so that generated commands are pre-filled on the next prompt line for
one-keystroke execution.
\end{abstract}

\begin{IEEEkeywords}
natural language interfaces, code generation, large language models,
parameter-efficient fine-tuning, quantization, LLM-as-judge, shell commands
\end{IEEEkeywords}

\section{Introduction}

The Unix shell is among the most powerful interfaces in computing, yet its
utility is gated by the user's ability to recall exact command syntax, flags,
and pipeline idioms. A user who knows \emph{what} they want---``find every
Python file changed in the last day and count the total lines''---may not recall
that this maps to
\texttt{find . -name '*.py' -mtime -1 -exec cat \{\} + | wc -l}.
Natural-language-to-command translation promises to remove this barrier, but a
practical assistant must satisfy three constraints simultaneously: it must
produce \emph{correct and complete} commands including the flags that encode
user intent, it must run \emph{locally} so that potentially sensitive requests
never leave the machine, and it must be \emph{fast and lightweight} enough to
feel native inside an interactive shell. Cloud-scale LLMs satisfy the first
constraint but violate the latter two.

This paper presents \textbf{nosh}, a natural-language-to-Bash assistant built on
a 1.5B-parameter model that runs entirely on consumer hardware. Our central
finding is methodological rather than architectural. A first fine-tuned model,
trained with reasonable hyperparameters on the corpora most readily available to
us---short curated command pairs and a large but tool-homogeneous
natural-language-to-Bash dataset dominated by \texttt{find}---taught the model to
comply with output formatting but did \emph{not} improve command correctness on
complex requests: on a held-out set of $600$ complex commands it scored $3.498$
against the base model's $3.503$ on a five-point scale, a difference within
noise. The model still hallucinated flags and arguments on anything beyond
simple commands. Diagnosing this, we found the training distribution never
contained the long, multi-tool pipelines the model was failing on.

Correcting the \emph{data} rather than the optimization recipe produced the
gain. Adding a corpus of complex, well-formed commands (average length $62$
characters, $28\%$ containing pipes) lifted the mean judge score on the same
held-out set to $4.178$, increased the fraction of commands rated $4$ or $5$
from $0.53$ to $0.74$, and nearly doubled exact-match accuracy from $0.26$ to
$0.46$. This supports a broader claim: for small code-generation models, when a
fine-tune fails to move a capability, the first suspect should be whether the
training distribution covers that capability's difficulty band, not the learning
rate or adapter rank.

Beyond quality, we treat deployability as a first-class research question. We
systematically study the trade-off between model footprint and generation
quality across four GGUF quantization tiers, and separately measure how many
in-context examples a fine-tuned model still needs. Both studies use the same
LLM-as-judge evaluation as our main experiments, so deployment decisions are
grounded in the same quality metric rather than proxy measures.

Our contributions are:
\begin{itemize}
  \item A demonstration, with controlled before/after evaluation, that
        training-data \emph{difficulty coverage}---not fine-tuning
        hyperparameters---is the binding constraint on complex-command quality
        for a 1.5B model, closing a $+0.68$ mean-judge-score gap on held-out
        complex commands.
  \item A reference-aware LLM-as-judge rubric for command generation that
        assigns partial credit across tool, flags, arguments, constraints, and
        safety, paired with deterministic syntax and formatting pre-checks.
  \item A deployment study over quantization tiers and few-shot lengths that
        identifies a $940$\,MB, two-shot configuration running in
        $\approx 1.7$\,GB of RAM at $1.4$--$1.6$\,s latency on CPU, integrated
        into the shell for one-keystroke execution.
\end{itemize}

\section{Related Work}

\textbf{Natural language to shell commands.} Translating English to shell
commands was framed as a semantic-parsing and neural machine-translation problem
by the NL2Bash effort of Lin \emph{et al.}~\cite{lin2018nl2bash}, which released
a corpus of English-to-Bash pairs and established token- and template-level
metrics. Earlier interactive systems such as Tellina explored neural translation
for file-system operations. These systems predate instruction-tuned LLMs and
generally target a narrower command vocabulary than modern code models; we reuse
the NL2Bash data as one component of our training mixture but find, consistent
with our central result, that its heavy skew toward \texttt{find} limits its
value for complex multi-tool commands.

\textbf{Code language models.} General code models such as Code
Llama~\cite{roziere2023codellama} and the Qwen2.5-Coder
series~\cite{hui2024qwen25coder} are pre-trained on large code corpora and
exhibit strong zero-shot shell competence, but at parameter counts and memory
footprints that preclude comfortable local, CPU-only deployment. We build on the
smallest instruction-tuned Qwen2.5-Coder checkpoint precisely to keep the
deployed footprint within a laptop's means.

\textbf{Parameter-efficient fine-tuning.} Low-Rank Adaptation
(LoRA)~\cite{hu2022lora} freezes the base weights and learns low-rank update
matrices, reducing the trainable-parameter count by orders of magnitude and
making single-GPU fine-tuning of billion-parameter models tractable. We use LoRA
throughout and merge the adapter into the base weights before quantization.

\textbf{Quantization for local inference.} Post-training quantization to 3--5
bits via the GGUF/\texttt{llama.cpp} $k$-quant formats~\cite{llamacpp} enables
CPU inference of billion-parameter models within gigabyte-scale memory budgets.
Prior work reports quality/size trade-offs at the level of perplexity; we instead
measure the trade-off directly in task quality using our command-generation
judge.

\textbf{LLM-as-judge evaluation.} Using a strong LLM to score the outputs of
another model correlates well with human preference and scales to large
evaluation sets~\cite{zheng2023judge}. Exact-match against a single reference
badly underestimates command-generation quality because many distinct commands
are functionally equivalent (e.g. \texttt{ls -la} vs. \texttt{ls -al}). We adopt
a reference-aware judge that treats the dataset command as a hint about intent
rather than the sole correct answer.

\section{System Design}

\subsection{Overview}

nosh comprises three stages: a data-preparation pipeline that assembles a
difficulty-balanced training mixture, a LoRA fine-tuning stage over
Qwen2.5-Coder-1.5B-Instruct, and a quantized local-inference command-line
application. Evaluation is performed by a separate LLM-as-judge stage that scores
generated commands against a multi-dimensional rubric. Fig.~\ref{fig:pipeline}
summarizes the flow.

\begin{figure}[t]
\centering
\footnotesize
\setlength{\fboxsep}{4pt}
\fbox{\parbox{0.92\columnwidth}{\centering
\textbf{Data pipeline} \\
curated pairs $+$ complex 6k corpus $+$ balanced NL2Bash \\
$\rightarrow$ comment stripping $\rightarrow$ leakage guard \\
$\downarrow$ \\
\textbf{Training} \\
Qwen2.5-Coder-1.5B $+$ LoRA (r{=}64) $\rightarrow$ merge 16-bit \\
$\downarrow$ \\
\textbf{Deployment} \\
GGUF quantize ($Q4\_K\_M$) $\rightarrow$ cold-load CLI $\rightarrow$ \texttt{print -z} \\
$\downarrow$ \\
\textbf{Evaluation} \\
deterministic checks $+$ Qwen3-14B reference-aware judge
}}
\caption{The nosh pipeline. Data assembly, parameter-efficient fine-tuning,
quantized deployment, and LLM-as-judge evaluation are decoupled stages.}
\label{fig:pipeline}
\end{figure}

\subsection{Data Preparation}

The training mixture is assembled from three sources with deliberately different
difficulty profiles. The first is a set of curated, domain-focused
natural-language/command pairs emphasizing developer and system-administration
tools (\texttt{git}, \texttt{docker}, \texttt{kubectl}, \texttt{sudo}). The
second is a $5{,}823$-example corpus of complex, well-formed commands whose
average length is $62$ characters and $28\%$ of which contain pipes---the
difficulty band the model most needed. The third is the NL2Bash
corpus~\cite{lin2018nl2bash}, which is large but heavily skewed: roughly $60\%$
of its commands invoke \texttt{find}. We counter this skew by capping the number
of examples per leading command token and sampling a balanced subset.

Two cleaning steps proved important. First, many complex-corpus entries append
explanatory comment lines (e.g. \texttt{\# [output of first 10 lines]}) beneath
the real command; training on these would teach the model to emit prose, in
direct violation of the ``output only a command'' contract, so we strip
pure-comment lines while preserving genuine \texttt{\#!} shebang scripts.
Second, we apply an explicit leakage guard that removes any training row whose
natural-language prompt appears in any evaluation set. The resulting training set
contains $9{,}086$ deduplicated examples spanning $890$ distinct leading command
tokens with a mean command length of $49.9$ characters. We hold out $600$
complex examples (\textsc{test-6k}) plus two curated $100$-example sets
(\textsc{009}, \textsc{010}) with zero prompt overlap into training.

\subsection{Model and Fine-Tuning}

We fine-tune Qwen2.5-Coder-1.5B-Instruct~\cite{hui2024qwen25coder} with
LoRA~\cite{hu2022lora} applied to all attention and MLP projection matrices. The
final model (V2) uses rank $64$ and scaling $\alpha=64$, trained for two epochs
with an effective batch size of $32$ (micro-batch $8$, gradient accumulation
$4$), a peak learning rate of $2\times10^{-4}$ under a cosine schedule, and
8-bit optimizer states. Training runs on a single 24\,GB consumer GPU. Loss
during in-training generation previews at full batch exceeded memory; disabling
previews and halving the micro-batch (with compensating accumulation) resolved
this without changing the effective batch size. After training, the LoRA adapter
is merged into the base weights in a separate process to produce a standalone
16-bit checkpoint suitable for quantization.

Prompts follow a fixed system instruction that constrains the model to emit only
a valid Bash command with no markdown, prose, or code fences, optionally
preceded by a small number of few-shot demonstrations. The same prompt shape is
used at training and inference time.

\subsection{Evaluation Rubric}

Because exact match against a single reference command underestimates quality, we
score generations with a reference-aware LLM-as-judge~\cite{zheng2023judge}. The
judge receives the instruction, the generated command, a reference command
(explicitly labeled a hint rather than ground truth), and two pre-computed
deterministic signals. It returns partial-credit scores along five dimensions:
tool appropriateness (boolean), flag coverage ($0$--$2$), argument fidelity
($0$--$2$), constraint satisfaction ($0$--$2$; e.g. ``read-only'', ``exclude
node\_modules'', a specific codec), safety (boolean; penalizing unrequested
destruction such as \texttt{rm -rf /}), and a holistic overall score
($1$--$5$). Two deterministic pre-checks feed the judge and are also reported
independently: a formatting check that rejects markdown, JSON, or conversational
preambles, and a syntax check that runs \texttt{bash -n} to confirm the command
parses.

\section{Experimental Setup}

\textbf{Models compared.} We compare three systems: the unmodified
Qwen2.5-Coder-1.5B-Instruct \emph{baseline}; \emph{V1}, a LoRA fine-tune trained
only on the curated pairs plus balanced NL2Bash; and \emph{V2}, the final
fine-tune that adds the complex-command corpus. V1 and V2 differ in both the
training mixture and the adapter rank ($32$ vs. $64$); the ablation of interest
is the mixture, and we discuss this confound in Section~\ref{sec:discussion}.

\textbf{Evaluation sets.} \textsc{test-6k} is the $600$-example held-out slice of
the complex corpus and is our primary benchmark. \textsc{009} and \textsc{010}
are $100$-example curated held-out sets used as secondary checks. No evaluation
prompt appears in training.

\textbf{Judge.} All judgments use Qwen3-14B (4-bit) as the scorer, run under
vLLM with greedy decoding for determinism. Deterministic checks (\texttt{bash
-n}, formatting) are computed in-process.

\textbf{Metrics.} We report the mean overall judge score
(\textsc{overall}, $1$--$5$), the fraction of commands rated at least $4$
(\textsc{pct}$\geq$\textsc{4}), normalized exact-match against the reference
(\textsc{exact}), and the \texttt{bash -n} pass rate (\textsc{syntax}).

\section{Results}

\subsection{Data Difficulty Drives Complex-Command Quality}

Table~\ref{tab:main} reports the three systems across all evaluation sets. On the
primary \textsc{test-6k} benchmark, V1 is statistically indistinguishable from
the untuned baseline ($3.498$ vs. $3.503$ overall), despite being a fine-tuned
model: it learned the output format---its formatting-compliance rate is
$0.998$---but not complex-command competence. V2, which differs from V1 chiefly
by the inclusion of the complex-command corpus, raises the overall score to
$4.178$, a $+0.68$ absolute gain, lifts \textsc{pct}$\geq$\textsc{4} from
$0.53$ to $0.74$, and nearly doubles exact-match from $0.26$ to $0.46$. The
per-dimension scores move consistently: argument fidelity rises from $1.45$ to
$1.71$ and constraint satisfaction from $1.26$ to $1.61$ (both out of $2$),
indicating the gain is concentrated exactly where V1 failed---in reproducing the
specific flags, values, and constraints of hard requests.

\begin{table}[t]
\caption{Main results across evaluation sets. \textsc{test-6k} is the primary
held-out complex benchmark ($n{=}600$); \textsc{009}/\textsc{010} are curated
held-out sets ($n{=}100$). Judge: Qwen3-14B. Best per column in bold.}
\label{tab:main}
\centering
\footnotesize
\begin{tabular}{llcccc}
\toprule
Set & Model & \textsc{overall} & \textsc{pct$\geq$4} & \textsc{exact} & \textsc{syntax} \\
\midrule
\multirow{3}{*}{\textsc{test-6k}}
 & Baseline & 3.503 & 0.533 & 0.257 & 0.973 \\
 & V1       & 3.498 & 0.548 & 0.230 & 0.968 \\
 & V2       & \textbf{4.178} & \textbf{0.740} & \textbf{0.455} & \textbf{0.990} \\
\midrule
\multirow{3}{*}{\textsc{009}}
 & Baseline & 3.640 & 0.560 & 0.230 & 0.920 \\
 & V1       & 3.830 & 0.630 & 0.290 & \textbf{1.000} \\
 & V2       & \textbf{3.990} & \textbf{0.680} & \textbf{0.320} & \textbf{1.000} \\
\midrule
\textsc{010} & V2 & 3.970 & 0.660 & 0.320 & 1.000 \\
\bottomrule
\end{tabular}
\end{table}

The curated \textsc{009}/\textsc{010} sets corroborate the trend at smaller
magnitude: V2 reaches $3.99$ and $3.97$ overall, improving on both baseline
($3.64$ on \textsc{009}) and V1 ($3.83$). The smaller gap here is expected, as
these sets contain shorter commands closer to V1's training distribution. The
\texttt{bash -n} syntax pass rate reaches $0.99$--$1.00$ for V2 across all sets,
confirming the model reliably emits parseable commands.

\subsection{Few-Shot Examples Are Largely Redundant After Fine-Tuning}

Because the deployed application re-encodes any few-shot prefix on every
invocation, we ask how many demonstrations the fine-tuned model still requires.
Table~\ref{tab:fewshot} varies the number of in-context example pairs from $0$
to $8$ at fixed quantization ($Q4\_K\_M$) on a $250$-example subset. The overall
score is flat within noise across all settings ($4.04$ to $4.10$), and even
zero-shot performance ($4.036$) is competitive. Fine-tuning has internalized the
task and the output contract, so few-shot prompting is nearly redundant. We
retain two examples in the default configuration for a marginal exact-match
benefit ($0.404$ vs. $0.388$ at zero-shot) and output-format anchoring at
negligible cost.

\begin{table}[t]
\caption{Few-shot length ablation ($Q4\_K\_M$, $n{=}250$). Scores are flat within
noise; the fine-tune has internalized the task.}
\label{tab:fewshot}
\centering
\footnotesize
\begin{tabular}{lccc}
\toprule
\# example pairs & \textsc{overall} & \textsc{pct$\geq$4} & \textsc{exact} \\
\midrule
0 & 4.036 & 0.716 & 0.388 \\
2 & 4.040 & 0.712 & 0.404 \\
4 & 4.104 & 0.716 & 0.404 \\
8 & 4.060 & 0.708 & 0.404 \\
\bottomrule
\end{tabular}
\end{table}

\subsection{Quantization Trade-off}

Table~\ref{tab:quant} evaluates four GGUF $k$-quant tiers at two-shot prompting
on a $100$-example subset. Quality is non-monotonic in bit-width: $Q4\_K\_M$
($940$\,MB) attains the highest overall score ($4.13$), while the larger
$Q5\_K\_M$ ($1073$\,MB) is no better ($3.99$), consistent with quantization
noise dominating the small remaining headroom. The more aggressive $Q3\_K\_M$
($786$\,MB) loses roughly $0.2$ overall and was additionally observed to emit
occasional non-UTF-8 bytes, a qualitative reliability concern. $Q4\_K\_M$ is thus
both the quality maximum and near the footprint minimum: no tier offers a
meaningful memory saving without a quality regression. We adopt $Q4\_K\_M$ as the
deployed configuration.

\begin{table}[t]
\caption{Quantization tier trade-off (two-shot, $n{=}100$). $Q4\_K\_M$ is the
quality maximum; larger tiers add memory without gain.}
\label{tab:quant}
\centering
\footnotesize
\begin{tabular}{lccccc}
\toprule
Tier & Size (MB) & \textsc{overall} & \textsc{pct$\geq$4} & \textsc{exact} & \textsc{syntax} \\
\midrule
$Q3\_K\_M$ & 786  & 3.94 & 0.67 & 0.33 & 0.99 \\
$Q4\_K\_S$ & 897  & 4.05 & 0.67 & 0.39 & 0.99 \\
$Q4\_K\_M$ & 940  & \textbf{4.13} & \textbf{0.72} & 0.39 & 0.97 \\
$Q5\_K\_M$ & 1073 & 3.99 & 0.70 & \textbf{0.40} & 0.99 \\
\bottomrule
\end{tabular}
\end{table}

\subsection{Deployment Characteristics}

The final application cold-loads the $940$\,MB quantized model on each
invocation---no persistent daemon---to keep the tool stateless and simple. On a
CPU-only configuration, representative of a laptop without a discrete GPU, end-to-end
latency including model load and generation is $1.4$--$1.6$\,s per query and peak
resident memory is $\approx 1.7$\,GB (model weights plus runtime and key/value
cache buffers at a $512$-token context). This footprint fits comfortably within
an $8$\,GB machine. The context window is capped at $512$ tokens, well above the
observed prompt-plus-command length, which bounds the key/value cache and reduces
compute. The generated command is emitted on standard output and injected into
the shell's next prompt line (via \texttt{print -z} in Zsh), so the user reviews
and executes it with a single keystroke.

\subsection{Qualitative Examples}

The deployed model produces idiomatic pipelines for requests well outside the
literal training strings. Representative outputs include:
\begin{lstlisting}
NL : show the top 10 most memory-hungry processes
CMD: ps aux --sort=-%mem | head -n 11

NL : count unique IP addresses in access.log
     sorted by frequency descending
CMD: awk '{print $1}' access.log | sort
     | uniq -c | sort -nr

NL : kill the process listening on port 8080
CMD: kill $(lsof -t -i:8080)
\end{lstlisting}
These exhibit correct multi-stage composition, appropriate flag selection, and
command substitution---capabilities absent from the V1 model.

\section{Discussion and Limitations}
\label{sec:discussion}

\textbf{Mixture/rank confound.} V1 and V2 differ in both training mixture and
adapter rank ($32$ vs. $64$). We attribute the gain to the mixture because the
improvement is concentrated on the complex held-out set and in the argument and
constraint dimensions that correspond directly to the added data, whereas rank
alone would be expected to help more uniformly. A rank-controlled ablation
(training rank $64$ on the V1 mixture) would isolate the effect fully and is left
to future work.

\textbf{Judge as oracle.} Our primary metric is a single LLM judge. While
reference-aware LLM judging correlates well with human
preference~\cite{zheng2023judge}, it can share blind spots with the generator
and does not execute commands. Sandbox execution-based evaluation would provide a
complementary, harder-to-game signal.

\textbf{Subset sizes in the deployment study.} The quantization and few-shot
sweeps use $100$- and $250$-example subsets; sub-$0.1$ score differences between
adjacent configurations are within sampling noise. The robust conclusions are
directional---``$Q4\_K\_M$ is the sweet spot and lower tiers regress,''
``few-shot examples are largely redundant''---rather than precise rankings among
statistically tied configurations.

\textbf{Safety.} The safety dimension penalizes unrequested destructive
operations at judging time, but nosh does not execute commands and inserts them
for user confirmation rather than auto-running them. Human-in-the-loop review
remains the primary safety mechanism.

\section{Conclusion}

We presented nosh, a 1.5B-parameter natural-language-to-Bash assistant that runs
on consumer hardware within $\approx 1.7$\,GB of memory. Our main finding is that
for small code-generation models, the difficulty coverage of the training
mixture---not the fine-tuning hyperparameters---was the binding constraint on
complex-command quality: adding a corpus of long, multi-tool commands raised the
mean LLM-judge score on held-out complex commands from $3.50$ to $4.18$ and
nearly doubled exact-match accuracy, whereas an earlier fine-tune lacking that
data matched the untuned baseline. We further showed, using the same
task-quality metric, that a fine-tuned model needs almost no few-shot examples
and that $4$-bit $Q4\_K\_M$ quantization is the quality/footprint sweet spot for
local deployment. Future work includes rank-controlled ablations to fully
isolate the data effect, execution-based evaluation in a sandbox, and extending
the difficulty-coverage principle to other structured-output generation tasks.

\begin{thebibliography}{00}
\bibitem{lin2018nl2bash} X. V. Lin, C. Wang, L. Zettlemoyer, and M. D. Ernst,
``NL2Bash: A Corpus and Semantic Parser for Natural Language Interface to the
Linux Operating System,'' in \emph{Proc. 11th Int. Conf. Language Resources and
Evaluation (LREC)}, 2018.
\bibitem{roziere2023codellama} B. Rozi\`ere \emph{et al.}, ``Code Llama: Open
Foundation Models for Code,'' \emph{arXiv preprint arXiv:2308.12950}, 2023.
\bibitem{hui2024qwen25coder} B. Hui \emph{et al.}, ``Qwen2.5-Coder Technical
Report,'' \emph{arXiv preprint arXiv:2409.12186}, 2024.
\bibitem{hu2022lora} E. J. Hu \emph{et al.}, ``LoRA: Low-Rank Adaptation of Large
Language Models,'' in \emph{Proc. Int. Conf. Learning Representations (ICLR)},
2022.
\bibitem{zheng2023judge} L. Zheng \emph{et al.}, ``Judging LLM-as-a-Judge with
MT-Bench and Chatbot Arena,'' in \emph{Advances in Neural Information Processing
Systems (NeurIPS)}, 2023.
\bibitem{llamacpp} G. Gerganov \emph{et al.}, ``llama.cpp: LLM inference in
C/C++,'' 2023. [Online]. Available: https://github.com/ggerganov/llama.cpp
\end{thebibliography}

\end{document}
